\documentclass{article}%
\usepackage[T1]{fontenc}%
\usepackage[utf8]{inputenc}%
\usepackage{lmodern}%
\usepackage{textcomp}%
\usepackage{lastpage}%
\usepackage{geometry}%
\geometry{margin=2cm}%
\usepackage{expex}%
\usepackage[T1]{fontenc}%
\usepackage[utf8]{inputenc}%
\usepackage[spanish]{babel}%
%
\title{Análisis Lingüístico: Gramatica Normativa Kaqchikel}%
\author{Generado por el TFM de Elías Carrasco}%
%
\begin{document}%
\normalsize%
\maketitle%
\section{Page 145}%
\label{sec:Page145}%
Demostrativo la.....la': %
\par\medskip%
\ex
\textit{Ese pájaro vuela mucho.} \\
\begingl
\gla Yalan nxik'an la jun tz'ikin la'. //
\glb yalan n-ø-xik'-a-n la jun tz'ikin la' //
\glc \textbf{INT} \textbf{INC-B3s-volar-BV-AP} \textbf{DEM} \textbf{IND} \textbf{pájaro} \textbf{DEM} //
\endgl
\xe
%
\medskip%
\par\medskip%
\ex
\textit{En ese pueblo hay árboles muy grandes.} \\
\begingl
\gla Pa la jun tinamit la' e k'o mama' taq che'. //
\glb Pa la jun tinamit la' e k'o mama' taq che' //
\glc \textbf{PRE} \textbf{DEM} \textbf{IND} \textbf{pueblo} \textbf{DEM} \textbf{B3p} \textbf{EST} \textbf{grande} \textbf{PL/DIM} \textbf{árbol} //
\endgl
\xe
%
\medskip%
De los anteriores demostrativos pueden acompañar al sustantivo solamente los antepuestos re y la, para indicar la visibilidad o cercanía del sustantivo. La partícula ri también funciona como demostrativo lo cual indica lejanía y sustantivo conocido. %
\subsection{Ejemplos:}%
\label{subsec:Ejemplos}%

%
Demostrativo re: %
\par\medskip%
\ex
\textit{...la señora que buscamos.} \\
\begingl
\gla ...re ixöq nqakanoj. //
\glb re ixöq n-ø-qa-kan-o-j //
\glc \textbf{DEM} \textbf{mujer} \textbf{INC-B3s-A1p-buscar-BV-SC} //
\endgl
\xe
%
\medskip%
\par\medskip%
\ex
\textit{Estos niños lo supieron.} \\
\begingl
\gla Re ak'wala' xketamaj. //
\glb Re ak'wal-a' x-ø-k-etam-a-j //
\glc \textbf{DEM} \textbf{niño-PL} \textbf{COM-B3s-A3p-saber-BV-VT} //
\endgl
\xe
%
\medskip%
Demostrativo la: %
\par\medskip%
\ex
\textit{...también ese árbol de pito...} \\
\begingl
\gla ...chuqa' la tz'ite' che'... //
\glb chuqa' la tz'ite' che' //
\glc \textbf{CONJ} \textbf{DEM} \textbf{palo de pito} \textbf{árbol} //
\endgl
\xe
%
\medskip

%
\end{document}